{\bfseries Transformada de Fourier clásica (DFT). Dada la señal discreta de longitud $N = 4$:
$\mathbf{x} = [x_0, x_1, x_2, x_3] = [1, 2, 0, -1]$.}

\begin{enumerate}[a)]
	\item {\bfseries Escriba la matriz de la DFT para $N = 4$ con $\omega_4 = e^{-2\pi i/4} = -i$.}

	\vspace{4cm}

	\item {\bfseries Calcule la transformada $\mathbf{X} = F\mathbf{x}$ componente por componente. Escriba el vector resultante $\mathbf{X} = [X_0, X_1, X_2, X_3]$.}

	\vspace{4cm}

	\item {\bfseries Verifique que se cumple la simetría conjugada para señales reales: $X_{4-k} = \overline{X_k}$.}

	\vspace{4cm}

	\item {\bfseries Dibuje aproximadamente la magnitud $|X_k|$ para $k = 0, 1, 2, 3$ y comente qué frecuencias están presentes.}

	\vspace{4cm}

\end{enumerate}
